\section{Conclusiones y propuesta de validación}
\label{sec:conclusiones}

En un año de rollos ligados entre el MC4 y el Skinpass de Pesquería, la corrección que el lazo de bending aplica sobre
la receta depende más del programa del Skinpass, o del periodo en que se procesa el rollo, que del producto: el
programa explica \ExcesoPrograma~puntos de su variación por encima del azar y el producto \ExcesoProducto. La corona con
que el rollo sale del MC4 se asocia con la corrección en la dirección física esperada, pero con una magnitud pequeña y
frágil, y la información previa al Skinpass no la anticipa claramente mejor que una receta por producto. Antes de
diferenciar recetas por producto, por lo tanto, proponemos validar con un experimento dentro de cada programa y
registrar la planeza de salida.

\paragraph{Contribución} Para Ternium, el artículo acota lo que una receta por producto podría absorber de la
corrección del Skinpass y muestra que la mayor parte de esa corrección cambia de un programa o de un día a otro; además
deja una unión verificada entre ambos molinos, reproducible con un solo comando. Para el campo, muestra cómo plantear
una pregunta respondible cuando la pregunta original exige un resultado que la planta no mide, y cómo someter una
asociación de planta a control negativo, multiverso y sensibilidad antes de traducirla en una decisión.

\subsection{Experimento propuesto}
\label{sec:experimento}
Proponemos una operación evolutiva~\cite{box1957}: perturbaciones pequeñas, dentro del rango ya operado, que no
detienen la producción.
\begin{itemize}
\item \emph{Qué se asigna.} Factor A: preset vigente frente a preset recentrado con la mediana histórica de la
      corrección de su celda producto $\times$ dimensión, limitado al rango de preset ya usado en esa celda. Factor B,
      sin costo adicional en un diseño factorial $2^2$: ajuste del preset por la corona relativa del MC4.
\item \emph{Unidad y bloqueo.} El rollo, alternado dentro de cada programa del Skinpass en bloques permutados de cuatro
      y estratificado por celda~\cite{box2005,montgomery2017}. Asignar por rollo y no por programa es la consecuencia
      práctica de H1: bloquear por programa elimina de la comparación la fuente de variación más grande. Es viable porque
      el preset ya cambia entre rollos de un mismo programa.
\item \emph{Resultado y criterio de decisión, fijados antes.} Resultado primario: $|\Y|$ promedio, en kN. Se adopta el
      recentrado si baja $|\Y|$ al menos \ExpDeltaA~kN con un intervalo que excluya cero y si ningún indicador
      secundario empeora más que un margen de no inferioridad acordado con Ternium~\cite{wellek2010}. Indicadores
      secundarios: defectos de inspector, rollos retenidos, elongación lograda, la cola (porcentaje de rollos con $|\Y|$
      mayor que su percentil 90 histórico) y, si se instala, la planeza de salida.
\item \emph{Volumen.} Con una desviación estándar de $|\Y|$ dentro de programa y celda de \ExpTodoSigmaDentro~kN, detectar
      \ExpDeltaA~kN con 80\,\% de potencia y $\alpha=0.05$ requiere \ExpTodoARollos{} rollos, unas \ExpTodoASemRollo{}
      semanas de producción del Skinpass (\ExpTodoRollosSemana{} rollos por semana). Asignar por programa requeriría
      \ExpTodoAProg{} programas (unas \ExpTodoASemProg{} semanas). Limitado a la familia más grande
      (\ExpFamilia, \ExpFamiliaRollos{} rollos al año), tomaría unas \ExpFamASemRollo{} semanas. Con ese volumen, el
      factor B solo se detectaría si su efecto fuera de al menos \ExpMDEB~kN, el doble del estimado.
\item \emph{Lo que el volumen revela.} Para probar lo que diferenciar por producto agrega al recentrado de la receta
      única (\ExpDeltaDif~kN) harían falta \ExpTodoDifRollos{} rollos, más de una década de producción. Ese es un
      hallazgo en sí mismo: con lo que hoy se mide, la diferenciación por producto no se puede validar
      experimentalmente en un plazo razonable, y no es donde está la mayor parte de la corrección.
\item \emph{Reglas de paro.} Revisión semanal por el ingeniero de proceso del Skinpass, que decide. Se detiene un
      brazo si su tasa de defectos de inspector supera a la del otro en más de un punto porcentual dos semanas seguidas, o
      si aparece un defecto de forma (fuera de planeza, ondulación, tableado) en rollos consecutivos del mismo brazo.
\end{itemize}

\subsection{Datos del Skinpass que harían respondible la pregunta original}
\begin{enumerate}
\item \emph{Planeza de salida} como medición continua a lo ancho, por rollo, con el tramo estable separado del enhebrado
      y la aceleración. Sin ella no hay un resultado de forma que comparar entre recetas.
\item \emph{Forma de entrada}, aunque sea una medición gruesa: es el principal confusor no medido del diagrama.
\item \emph{Consigna y valor medido de cada palanca} (bending, elongación, fuerza, tensiones, velocidad) a lo largo del
      rollo, y la \emph{lógica de los lazos}: qué actuador responde a qué medición. Sin esto no se distinguen las
      palancas de los mediadores.
\item \emph{Campaña de rodillos de trabajo} de cada rollo (fecha de montaje y toneladas acumuladas), el \emph{historial
      de recetas} con fechas de vigencia y los \emph{ajustes del operador} con su motivo.
\item \emph{El mismo identificador} que usa el MC4: con él ligamos el \PctCalienteEnSKP--\PctDecapadoEnSKP\,\% de los
      rollos planeados para el Skinpass. Las señales de 2026 que recibimos durante la Entrega~I no coincidían con ningún
      rollo de 2025, así que cualquier dato nuevo debe cubrir el mismo periodo que el registro del MC4.
\end{enumerate}
Una advertencia acompaña esta lista: mientras exista una receta por tipo de material, la variación de configuración
dentro de cada producto será pequeña y en buena parte reactiva. Los registros históricos, aun completos, medirán bien
cuánto difieren los productos en lo que llegan necesitando, pero la diferencia de respuesta, que es lo que justificaría
una receta por producto, casi con seguridad exigirá un experimento diseñado en el Skinpass.

En suma, es verdad que la corrección del Skinpass depende más del programa o del periodo que del producto, y poco de la
forma con que el rollo sale del MC4. Sigue sin saberse si reducir esa corrección mejora la planeza y si algún producto
responde distinto a la receta. Lo resolvería un experimento rollo a rollo dentro de cada programa, con la planeza de
salida medida.
